\section{Does the split measure what it claims?}\label{sec:exp}

A released checkpoint cannot validate the split, because what the model can see is unknown.
We validate it where the answer is known, first in simulation and then on Visual Genome
predictors whose information content we control. All intervals are 95\%.

\paragraph{Simulation.}\label{sec:simulation} With 24 cells over five labels, an old model
equal to the cell prior and a new one interpolating towards an oracle, the estimated
case-level gain grows linearly with the injected signal; under a within-cell null the
randomization test rejects in $0.030$ of 200 datasets at level $0.05$ (mean estimate
$-8.97\times10^{-5}$). For the interval we use on real data, we simulate relations
clustered into images with a shared per-image effect and Zipf-distributed cell sizes (most
cells hold two cases, most relations sit in a few large ones): the image-clustered interval
covers in $94.0\%$ of 500 runs, against $88.6\%$ for an interval that treats relations as
independent. With VG150's own cell sizes it covers $96.3\%$; only when every cell holds two
cases does it fall short, to $81.5\%$ (\cref{app:clt}).

Three designs check that the two parts separate what they claim to. A \emph{prior-only}
improvement, where the new model corrects a noisy estimate of the cell prior, gives a
case-level gain centred at zero ($0.0002$, sd $0.0035$). An \emph{unrecorded within-cell
shortcut} --- a binary feature varying inside each cell that the new model uses --- is
credited to the case-level part ($+0.041$, sd $0.004$): the split measures all within-cell
signal relative to the declared grouping, shortcut or not. And a \emph{calibration-only}
change, where the new model is a temperature-softened copy of an overconfident one carrying
identical case-level information, is credited a large spurious case-level \emph{loss} of
$-0.488$ on raw outputs, which the confidence-matching protocol of \cref{sec:matching}
reduces to $-0.0003$ (sd $0.0017$). That check is what matching is built to pass. A
\emph{class-wise shift of the operating point}, adding a fixed logit offset per label to the
old model, leaves every within-cell ranking unchanged, yet survives temperature matching
($+0.0106$ for a logit-adjustment offset) and vanishes only when each model is matched by a
temperature and a per-class bias ($0.0000$, sd $0.0003$; \cref{app:rankrob}). Both
protocols are therefore part of the procedure.

\paragraph{Controlled predictors on VG150.}\label{sec:data} We train predicate classifiers whose inputs we control on a VG150-style split we
reconstruct from the released Visual Genome annotations, with the same 150-object and
50-predicate vocabulary and images divided 70/30 between fitting and audit
(\cref{app:vgmain-full}); it is not the canonical split of \cref{sec:sggaudit}. The predictors are: a frequency table that sees only the class
pair (FREQ), an MLP on class embeddings (MLP-CLASS), and MLPs adding twelve box-geometry
features (MLP-SPATIAL, MLP-SPATIAL+). The audit covers $227{,}337$ of $229{,}605$ test
relations ($99.0\%$) in $6{,}346$ class-pair cells.

MLP-CLASS beats FREQ by $0.03386$ of quadratic score and the split assigns \emph{all} of it
to the group-level part: the case-level gain is $0.00000$ with a zero-width interval. That zero is algebraic, not empirical --- both models are
functions of the class pair alone, so \cref{thm:population} forces it --- and it is the
sharpest check that nothing is credited to case-level evidence when none exists. Adding box
geometry then yields a case-level gain of $0.01439$ (CI $[0.01373,0.01504]$, $p=.002$) that no
aggregate score isolates; against MLP-CLASS the total gain of $0.00885$ even
\emph{understates} it, because the group-level part falls by $0.00554$.
